\begin{thebibliography}{99}

\bibitem{Project3} FSPT project draft, \emph{Lattice Model of 3+1D Interacting Fermionic Symmetry-Protected Topological Phases: Complete Classification, Stacking Group and Space Groups}, supplied file \texttt{Stacking\_Group\_of\_3\_1D\_FSPT (4).pdf}, dated 10 October 2026. In particular Secs.~VI~B--F and Appendices A, E, G, H.

\bibitem{Project4} FSPT project draft, \emph{Lattice Model of 4+1D Interacting Fermionic Symmetry-Protected Topological Phases: Complete Classification and AHSS}, supplied file \texttt{Lattice\_Model\_of\_4\_1D\_FSPT (11).pdf}, dated 10 October 2026. In particular Sec.~V and Appendices A, F, G.

\bibitem{WQFCG} Q.-R. Wang, Y. Qi, C. Fang, M. Cheng, and Z.-C. Gu, \emph{Exactly solvable lattice models for interacting electronic insulators in two dimensions}, Phys. Rev. B \textbf{108}, L121104 (2023), arXiv:2112.15533.

\bibitem{BCHM} M. Barkeshli, Y.-A. Chen, P.-S. Hsin, and N. Manjunath, \emph{Classification of (2+1)D invertible fermionic topological phases with symmetry}, Phys. Rev. B \textbf{105}, 235143 (2022), arXiv:2109.11039.

\bibitem{BrumfielMorgan} G. Brumfiel and J. Morgan, \emph{The Pontrjagin Dual of 4-Dimensional Spin Bordism}, arXiv:1803.08147v1 (2018), especially Secs.~1.5, 6--8. Their physical dimension and coefficient tower differ from the present $4+1$D problem.

\bibitem{SteenrodEpstein} N. E. Steenrod, lectures written and revised by D. B. A. Epstein, \emph{Cohomology Operations}, Annals of Mathematics Studies \textbf{50}, Princeton University Press (1962), especially Chapters I and VI--VIII.

\bibitem{MosherTangora} R. E. Mosher and M. C. Tangora, \emph{Cohomology Operations and Applications in Homotopy Theory}, Harper \& Row (1968), especially the construction of operations, spectral sequences, and Postnikov systems.

\bibitem{FreedHopkins} D. S. Freed and M. J. Hopkins, \emph{Reflection positivity and invertible topological phases}, Geometry \& Topology \textbf{25} (2021), 1165--1330, arXiv:1604.06527v6 (2022).

\bibitem{SerreEM} J.-P. Serre, \emph{Cohomologie modulo 2 des complexes d'Eilenberg--MacLane}, Commentarii Mathematici Helvetici \textbf{27} (1953), 198--232, doi:10.1007/BF02564562.

\bibitem{HatcherSS} A. Hatcher, \emph{Spectral Sequences}, preliminary additional chapter of \emph{Algebraic Topology}, author manuscript, \url{https://pi.math.cornell.edu/~hatcher/AT/ATch5.pdf}; especially the Serre sequence and Eilenberg--Mac Lane space calculations.

\bibitem{McClearyBSS} J. McCleary, \emph{A User's Guide to Spectral Sequences}, 2nd ed., Cambridge Studies in Advanced Mathematics \textbf{58}, Cambridge University Press (2001), Chapter 10, \emph{The Bockstein Spectral Sequence}. Author-hosted chapter: \url{https://pages.vassar.edu/mccleary/files/2011/04/MC10.fin_.pdf}.

\bibitem{DebrayYeYu} A. Debray, W. Ye and M. Yu, \emph{Bosonization and Anomaly Indicators of (2+1)-D Fermionic Topological Orders}, arXiv:2312.13341v2 (2025), especially Sec.~6.2 and Proposition~6.27. We use the specified tangential structures and generating manifolds, not a crystalline classification table.

\bibitem{AHSSGuide} Phy-Ren, \emph{FSPT\_AHSS}, insulator branch, \emph{Formula Guide}, dimensional formulas, operation conventions and exact implementation. Public formula guide rechecked 16 October 2026. \url{https://github.com/Phy-Ren/FSPT_AHSS/tree/insulator}.

\bibitem{STZ} S. Saito, Y. Tachikawa and Y. Zhang, \emph{Bosonic SPT and invertible phases and its relation to Steenrod's problem}, arXiv:2607.14662v2, Sections 2.2 and 3.2.3. Used for the oriented free response lattice, not as an electronic classification table.

\bibitem{AW} B. Antieau and B. Williams, \emph{The topological period--index problem over 6-complexes}, arXiv:1208.4430v2, Theorem A and Section 5. Its target is twisted $K$-theory; no interacting-spectrum equivalence is imported.

\bibitem{repo_d3} Same repository, \emph{Integral first obstruction for the optional intrinsic 4+1D row},
\href{https://github.com/Phy-Ren/FSPT_AHSS/blob/insulator/docs/e8/research/parity_bridge/intrinsic_d3_unitary_4p1.md}{\texttt{intrinsic\_d3\_unitary\_4p1.md}}, reviewed 15 October 2026.

\bibitem{repo_model} Same repository, \emph{Physical meaning of the optional intrinsic coefficient in 4+1D},
\href{https://github.com/Phy-Ren/FSPT_AHSS/blob/insulator/docs/e8/research/parity_bridge/intrinsic_charge3_realization_4p1.md}{\texttt{intrinsic\_charge3\_realization\_4p1.md}}, reviewed 15 October 2026.

\bibitem{HanLi} F. Han and Y. Li,
\emph{Differential Models for the Anderson Dual to Twisted Spin-c-Bordism and a Twisted Anomaly Map},
\href{https://arxiv.org/abs/2510.27286v2}{arXiv:2510.27286v2},
Sections 3.1, 3.4 and 4.3. The matrix in this note is not attributed as a printed
formula of that paper.

\bibitem{CLV} X.~Chen, Y.-M.~Lu and A.~Vishwanath, \emph{Symmetry protected topological phases from decorated domain walls}, arXiv:1303.4301.

\bibitem{SXH} H.~Song, C.~Z.~Xiong and S.-J.~Huang, \emph{Bosonic Crystalline Symmetry Protected Topological Phases Beyond the Group Cohomology Proposal}, arXiv:1811.06558.

\bibitem{YC} X.~Yang and M.~Cheng, \emph{Gapped boundary of $(4+1)$d beyond-cohomology bosonic SPT phase}, Phys. Rev. B \textbf{110}, 045137 (2024), arXiv:2303.00719v3, especially Section II and Appendix B.

\bibitem{RealSpin} H. H. Abdallah and Y. Kamel, \emph{The homotopy fixed points of Real spin bordism}, arXiv:2511.10624 (2025). The equivariant splitting is not claimed to identify the gerbe-dependent line action used here.

\bibitem{Metlitski} M. A. Metlitski, \emph{$S$-duality of $u(1)$ gauge theory with $\theta=\pi$ on non-orientable manifolds: Applications to topological insulators and superconductors}, arXiv:1510.05663.

\bibitem{SmithFiber} A. Debray, S. K. Devalapurkar, C. Krulewski, Y. L. Liu, N. Pacheco-Tallaj and R. Thorngren, \emph{The Smith Fiber Sequence and Invertible Field Theories}, arXiv:2405.04649, especially Theorem~5.1, Corollary~5.8 and Example~7.11. The paper proves the general exact sequence; its application to $H_-$ in this note is derived explicitly.

\bibitem{QHZ} X.-L. Qi, T. L. Hughes and S.-C. Zhang,
\emph{Topological Field Theory of Time-Reversal Invariant Insulators},
\href{https://arxiv.org/abs/0802.3537}{arXiv:0802.3537}, Section III.

\bibitem{BahriGilkey} A. Bahri and P. Gilkey, \emph{The eta invariant, Pin$^c$ bordism, and equivariant Spin$^c$ bordism for cyclic 2-groups}, Pacific J. Math. \textbf{128} (1987), 1--24, doi:10.2140/pjm.1987.128.1. Theorem 0.1 is used only for powers of two. The fixed-point eta formula and coefficient input are separately checked in the new applications.
\bibitem{Hsieh} C.-T. Hsieh, \emph{Discrete gauge anomalies revisited}, arXiv:1808.02881v1 (2018), especially Eqs. (2.35), (2.43), (2.49), and footnote 7. The corrected augmentation-ideal exponent in that footnote is retained.

\end{thebibliography}
